\section{Conclusion}
\label{sec:conclusion}

We have introduced ApportionRegions (AR), an algorithm that
generates congressional district maps by applying the prime factorization of
the apportioned seat count as a hierarchical geographic partition. The
algorithm has two inputs: arithmetic (the prime factorization of $n$) and
geography (the census-tract adjacency graph with TIGER boundary weights).
It produces a unique map for each state at each seat count, with no political
inputs and no human boundary choices.

The key structural properties --- reuse of canonical prime splits across
seat counts, stability under reapportionment when factorizations share a
common prefix, and complete structural breaks when they do not --- follow
directly from the Fundamental Theorem of Arithmetic. The map is as
non-arbitrary as the seat count itself.

The constitutional argument is that AR extends Huntington-Hill
(Article~I~\S2) from seat allocation to geographic allocation, rather than
invoking the Manner clause (\S4) as a separate authority. The apportionment
process, taken to its geographic conclusion, produces the map.

The most striking empirical finding is the $51 \to 52$ California transition,
where consecutive seat counts have completely disjoint prime structures,
forcing a full map redesign. This is not a defect of AR --- it is a
consequence of the mathematics. Any redistricting algorithm that takes the
arithmetic of seat counts seriously will face the same transitions. AR makes
them visible and principled.

Future work: seed sensitivity of $k$-way METIS for large prime $k$;
formal proof of uniqueness for prime-$k$ minimum cuts; empirical
disruption study for 2010$\to$2020 reapportionment transitions;
and the AR-smooth variant for minimising tract reassignments.
